\subsection{El tamaño de la historia domina el error de valuación}

El patrón más fuerte es la reducción del error al aumentar la cantidad de información histórica. La Tabla \ref{tab:rmse_n} presenta el RMSE promedio sobre las 27 combinaciones de $(\sigma_0,K/S_0,T)$ disponibles para cada tamaño de muestra.

\begin{table}[H]
\centering
\caption{RMSE promedio por tamaño de historia y regla de valuación}
\label{tab:rmse_n}
\begin{tabular}{rcccc}
\toprule
$n$ & Full Bayes & Media posterior & MAP & MLE \\
\midrule
63   & 0.38031 & 0.37972 & 0.38928 & 0.38066 \\
252  & 0.18289 & 0.18282 & 0.18692 & 0.18288 \\
1260 & 0.08416 & 0.08415 & 0.08595 & 0.08413 \\
\bottomrule
\end{tabular}

\vspace{0.15cm}
\footnotesize Nota: promedio sobre 27 regímenes contractuales y de volatilidad por fila; 2,000 muestras históricas independientes por escenario.
\end{table}

Para full Bayes, la mediana del cociente $\RMSE_{63}/\RMSE_{1260}$ a través de contratos y volatilidades es 4.67. Los cocientes medianos correspondientes son 4.66 para media posterior, 4.50 para MAP y 4.55 para MLE. Por tanto, pasar de aproximadamente tres meses a cinco años de historia diaria reduce el error en un orden cuantitativamente mucho mayor que cambiar entre full Bayes, media posterior y MLE.

La Tabla \ref{tab:rmse_n} también muestra que no existe un ganador universal. La media posterior presenta el menor RMSE promedio en $n=63$ y $n=252$, mientras MLE lo hace en $n=1260$, pero las diferencias entre estas tres reglas son pequeñas. MAP es la única regla con una penalización promedio persistente.

Una comparación escenario por escenario produce un matiz adicional. Full Bayes obtiene el menor RMSE en 35 de 81 escenarios, MLE en 35 y MAP en 11; la media posterior no es el mínimo exacto en ningún escenario. Este conteo no contradice sus bajos promedios: cuando las diferencias son diminutas, una regla puede ser sistemáticamente segunda por un margen mínimo y aun así exhibir un promedio competitivo. Por ello, los conteos de ``ganadores'' se interpretan sólo como complemento de las magnitudes de error.

\subsection{Full Bayes frente al plug-in de media posterior}

La comparación central del artículo es entre \eqref{eq:fb} y \eqref{eq:pm}. Full Bayes presenta menor RMSE en 45 de 81 escenarios y el plug-in en 36. La diferencia mediana $\RMSE_{FB}-\RMSE_{PM}$ es $-1.26\times10^{-6}$, mientras que la brecha absoluta relativa mediana es apenas 0.045\%. En términos de punto de precio, la mediana de $|\Delta_{FB,PM}|$ es 0.000736 y el máximo observado en el diseño completo es 0.01896 unidades monetarias.

Estas cifras confirman que la integración posterior completa rara vez cambia materialmente el punto de valuación dentro de la especificación Black--Scholes estudiada. Sin embargo, el experimento factorial permite identificar con mayor precisión cuándo la diferencia deja de ser despreciable. La Tabla \ref{tab:fbpm_n} muestra que la brecha cae rápidamente con la cantidad de información histórica.

\begin{table}[H]
\centering
\caption{Brecha Full Bayes--media posterior según información histórica}
\label{tab:fbpm_n}
\begin{tabular}{rccc}
\toprule
$n$ & Mediana $|\Delta_{FB,PM}|$ & Máximo $|\Delta_{FB,PM}|$ & Mediana brecha RMSE relativa \\
\midrule
63   & 0.007927 & 0.018964 & 0.288\% \\
252  & 0.001962 & 0.004788 & 0.045\% \\
1260 & 0.000391 & 0.000955 & 0.007\% \\
\bottomrule
\end{tabular}

\vspace{0.15cm}
\footnotesize La brecha relativa se calcula como $|\RMSE_{FB}-\RMSE_{PM}|/\RMSE_{PM}$.
\end{table}

El patrón es consistente con \eqref{eq:taylor}: al concentrarse $p(\sigma\mid\mathcal D)$, el término proporcional a $\Var(\sigma\mid\mathcal D)$ disminuye. La reducción no depende de que el payoff deje de ser dependiente de trayectoria; surge de que la incertidumbre de volatilidad que alimenta ese payoff se vuelve pequeña.

\subsection{MAP, MLE y ausencia de dominancia universal}

Full Bayes tiene menor RMSE que MAP en 60 de 81 escenarios. En promedio, $\RMSE_{FB}-\RMSE_{MAP}=-0.00493$, y la diferencia mediana es $-0.00232$. La debilidad del MAP se observa también en sesgo: con $n=63$, su sesgo absoluto promedio a través de regímenes es 0.1046, frente a 0.0329 de full Bayes y 0.0323 del plug-in de media posterior. Con $n=252$ los valores correspondientes son 0.0259, 0.0065 y 0.0067; con $n=1260$, 0.00434, 0.00135 y 0.00111. En esta especificación, la moda posterior no constituye una regla de pricing especialmente robusta.

La comparación con MLE es más equilibrada. MLE presenta menor RMSE en 46 escenarios frente a 35 para full Bayes, pero la diferencia mediana $\RMSE_{FB}-\RMSE_{MLE}$ es sólo $2.4\times10^{-5}$ y la diferencia media es $-1.0\times10^{-4}$. Con historias largas, MLE obtiene el mínimo en 20 de 27 escenarios de $n=1260$, lo cual es compatible con la convergencia entre estimadores cuando la verosimilitud domina y la posterior se concentra.

Estos resultados no respaldan una narrativa de superioridad universal del enfoque bayesiano en error cuadrático. Su ventaja distintiva es otra: full Bayes produce simultáneamente un punto de integración y una distribución de precios inducida por incertidumbre paramétrica. Cuando la posterior es estrecha, un plug-in o MLE puede reproducir el punto de precio con precisión similar; cuando la posterior es más dispersa y el funcional es curvo, la integración empieza a separarse de las reglas puntuales.
